\chapter{Discussion}\label{ch:discussion}

\section{Introduction}
This chapter examines what the results mean, what they do not show, and what was learned building the system. It discusses task-dependent baseline results and why short-chunk policies like GR00T~N1.7 run dry whatever the client does, alongside what adaptive playback can and cannot do. It then discusses what the faults cost on the robot, what \drtc{} wins back and what the stack adds. It concludes with four implementation lessons and limitations including sample size, shared servers and simulated faults.

\section{What the baselines show}
The baselines (Section~\ref{sec:baselines}) show, first, that size and recency do not order the policies. GR00T~N1.7
and \smolvla{} were the strongest overall, at 90\% on sorting and 70\% and 60\% on cube in box, yet \smolvla{} is one of the smallest
models (0.45B parameters); Diffusion Policy, which has no language input, was the best on cube in box (90\%); and
\groot{} was weak on both tasks (20\% and 10\%) while its successor N1.7, trained on the same data, was far better.
Second, the ranking depends on the task: \pizerofive{} did better on cube in box (60\%) than on sorting (40\%), where the
policy must read the cube's colour and the layout, and \smolvla{} the reverse. Third, ten trials per policy and task
give wide intervals (Tables~\ref{tab:basesort} and~\ref{tab:basecube}), so only large gaps, such as N1.7 against N1.5,
are clear. MolmoAct~2 never succeeded; we treat its LoRA fine-tune as broken rather than as a measure of the model. No
baseline trial was ended by a fault, so every failure was the policy's own, which makes the baselines the reference
against which the faults of E4 are measured.

\section{\drtc{} on the robot}
The robot experiment answers the problem statement for one task and four policies. \textbf{What the faults cost:} a
fault mix of about 3.5~s in a 30~s trial lowered the patient client's success from 68 to 13\%, not because the policies
failed but because the client did: after a lost reply it sends one fresh request, and when that is lost too it waits
until LeRobot's start-up deadline ends the run (32 of 40 faulted runs). \textbf{What \drtc{} wins back:} all of it,
within the intervals (65\%), because its cooldown replaces a lost reply with a request carrying a fresh observation;
the arm moved again after the outage in 31 of 40 trials, about a second after the link returned. The bench predicted
this (E1: \drtc{} recovered from every fault, the patient client from about a fifth
to a half), and E2 located it in the
cooldown. \textbf{What the stack adds:} smoother motion, about three times lower jerk, at no cost in success, and fewer
stalls for the short-chunk policies through adaptive playback; the price is a deviation of about 1--2\degree{} from the
policy's own actions and, for \pizerofive, a slower recovery after the latency spike.

Two qualifications matter. With ten trials per policy, only one per-policy comparison survives the pre-registered Holm
correction (GR00T~N1.7, \drtc{} against faults); the pooled effects are large but exploratory. And the benefit depends
on the policy. \pizerofive{}, whose 50-action chunk outlasts short faults, lost least under faults (8 to 4) and gained
least from \drtc{} (4 to 6), as Figure~\ref{fig:whendrtc} predicts; the short-chunk GR00T~N1.7 and Diffusion Policy
lost almost everything and gained the most. \drtc{} matters most where a policy cannot ride out a fault on its own
queue.

\section{Short-chunk policies and the round trip}\label{sec:shortchunk}
A policy can only keep the arm moving if a new chunk arrives before the current one is used up, so the chunk must last
longer than the round trip. Figure~\ref{fig:whendrtc} places the policies by their chunk duration and their median
turnaround from the Pi. ACT, \smolvla{} and \pizerofive{} lie well above the diagonal: on a good link any client keeps
their queue full, and a scheduler matters when a reply is late or lost. The two GR00T models lie below it: their
16-action chunk lasts 0.53~s, while their replies took 0.86 and 0.97~s at the median, so the queue runs dry before
every reply, whatever the client does. For such policies the remedy is not the scheduler but the playback: adaptive
playback slows the arm when few actions are left, stretching each chunk at the cost of speed, and a shorter round trip
(the cable, compressed images) moves the policy towards the diagonal. Diffusion Policy, at 0.8~s for 0.8~s, sits on
the line.

\diagram{whendrtc}{When \drtc{} can help: chunk length against round trip.}{fig:whendrtc}


\section{Lessons from the implementation}

\begin{itemize}
\item \textbf{The slow link finds the bugs.} Three of the four fixes of Table~\ref{tab:fixes} concern the cooldown, and
      all three appeared only with a slow link and a short chunk. The replay bench with the latency measured on the
      robot reproduced them in minutes and became the main development tool.
\item \textbf{A method can hurt in combination.} The soft stop looked harmless alone but, by keeping the queue alive, let
      a reply from an old observation be accepted, and the next chunk jumped; it was dropped
      (Section~\ref{sec:smoothing}). Each smoothing method is therefore tested alone before the stack is chosen.
\item \textbf{Measure what the system does, not what the operator presses.} In an early run several of the
      trials that the client ended were stored as ordinary failures, because the operator pressed the failure key
      while the stopped arm returned home. The label is now read from the run's own record.
\item \textbf{Fix the analysis before the data.} The comparisons that are tested and the rule that chooses the stack
      were written down before the bench ran, so the result cannot shape the question.
\end{itemize}

\section{Limitations}\label{sec:limits}

\begin{itemize}
\item \textbf{Sample size.} Ten trials per cell give wide intervals (a 50\% success rate over ten trials has a 95\%
      interval of about 24--76\%), so only large differences in success are significant. The robot plans can grow to
      twenty trials.
\item \textbf{Scope on the robot.} The robot \drtc{} experiment covers one task (cube in box) and four policies; the bench covers two
      tasks and seven policies, without task success.
\item \textbf{Operator judgement.} Success is judged by the operator, who knows the condition; the success rule is
      written down before the first trial to limit this.
\item \textbf{Shared servers.} Other users' jobs run on the A100 and H200, so inference times include contention.
\item \textbf{A slower link.} \drtc's cooldown is set from the replies seen so far, so a reply that is only slower
      than the recent ones can be given up as lost; over the Pi's WiFi this abandoned every reply of the slower
      policies (Section~\ref{sec:edge-cable}). A cooldown that adapts faster to a slower link is future work.
\item \textbf{Simulated faults.} The faults follow fixed shapes, shifted per trial, rather than a recorded real network,
      and the bench simulates the latency.
\item \textbf{Uncompressed images.} About 260~ms of each round trip is the network share, mostly the 2.9~MB of raw
      images per request; compression would shorten it and change the timing.
\end{itemize}

\section{Summary}
This chapter discussed what the baselines show: neither size nor recency orders the policies, the ranking depends on
the task, and only large gaps are clear with ten trials. It explained why the short-chunk GR00T policies run dry
whatever the client does, since their chunk is shorter than their round trip from the Pi, and why adaptive playback
or a shorter round trip, not the scheduler, is the remedy there. On the robot, the faults cost the patient client most of its
successes, \drtc{} recovered them through its cooldown, and the stack made the motion smoother without lowering
success. It collected the lessons of building the system and its limitations.
